\documentclass[../../../include/open-logic-section]{subfiles}

\begin{document}

\olfileid{sth}{spine}{Valphabasic}
\olsection{ویژگی‌های پایهٔ مراحل}

برای آشکارکردن اهمیت بنیادیِ تعریف~$V_\alpha$ها، چند نتیجهٔ پایه دربارهٔ
آن‌ها ارائه می‌کنیم. با یک تعریف آغاز می‌کنیم:\footnote{برای «توانمند»
اصطلاح معیار واحدی وجود ندارد؛ این نام را~\citet{ButtonLT1} به کار
می‌برد.}
\begin{defn}
	مجموعهٔ~$A$ \emph{توانمند} است اگر و تنها اگر
	$\forall x((\exists y \in A)x \subseteq y \lif x \in A)$.
\end{defn}
\begin{lem}\ollabel{Valphabasicprops}
برای هر عدد ترتیبی~$\alpha$:
\begin{enumerate}
	\item\ollabel{Valphatrans} $V_\alpha$ متعدی است.
	\item\ollabel{Valphapotent} $V_\alpha$ توانمند است.
	\item\ollabel{Valphacum} اگر~$\gamma\in\alpha$، آنگاه
	$V_\gamma\in V_\alpha$ (و در نتیجه بنا بر~\olref{Valphatrans}،
	$V_\gamma\subseteq V_\alpha$ نیز برقرار است).
\end{enumerate}
\end{lem}

\begin{proof}
این سه ادعا را هم‌زمان با استقرای ترامتناهی اثبات می‌کنیم. برای فرض
استقرا، فرض کنید~\olref{Valphatrans}--\olref{Valphacum} برای هر عدد
ترتیبی~$\beta<\alpha$ برقرار باشند.

حالت~$\alpha=\emptyset$ بدیهی است.

فرض کنید~$\alpha=\ordsucc{\beta}$. برای اثبات~\olref{Valphacum}، اگر
$\gamma\in\alpha$، آنگاه یا~$\gamma\in\beta$ یا~$\gamma=\beta$.
در حالت نخست بنا بر فرض استقرا و در حالت دوم بنا بر بازتابی‌بودن شمول،
$V_\gamma\subseteq V_\beta$؛ پس
$V_\gamma\in\Pow{V_\beta}=V_\alpha$. برای اثبات
\olref{Valphapotent}، فرض کنید
$A\subseteq B\in V_\alpha$؛ یعنی
$A\subseteq B\subseteq V_\beta$. در نتیجه~$A\subseteq V_\beta$ و پس
$A\in V_\alpha$. برای اثبات~\olref{Valphatrans}، توجه کنید اگر
$x\in A\in V_\alpha$، آنگاه~$A\subseteq V_\beta$، پس~$x\in V_\beta$.
چون بنا بر فرض استقرا~$V_\beta$ متعدی است، $x\subseteq V_\beta$؛ و در
نتیجه~$x\in V_\alpha$.

فرض کنید~$\alpha$ عدد ترتیبی حدی باشد. برای اثبات~\olref{Valphacum}،
اگر~$\gamma\in\alpha$، آنگاه
$\gamma\in\ordsucc{\gamma}\in\alpha$؛ پس بنا بر فرض استقرا
$V_\gamma\in V_{\ordsucc{\gamma}}$ و بنابراین
$V_\gamma\in\bigcup_{\beta\in\alpha}V_\beta=V_\alpha$. برای اثبات
\olref{Valphatrans} و~\olref{Valphapotent} کافی است توجه کنیم اجتماع
مجموعه‌های متعدی (و به‌ترتیب، توانمند) متعدی (و به‌ترتیب، توانمند) است.
\end{proof}

\begin{lem}\ollabel{Valphanotref}
برای هر عدد ترتیبی~$\alpha$، $V_\alpha\notin V_\alpha$.
\end{lem}

\begin{proof}
با استقرای ترامتناهی. آشکار است که
$V_\emptyset\notin V_\emptyset$.

اگر~$V_{\ordsucc{\alpha}}\in V_{\ordsucc{\alpha}}=
\Pow{V_\alpha}$، آنگاه
$V_{\ordsucc{\alpha}}\subseteq V_\alpha$؛ و چون بنا بر
\olref{Valphabasicprops}، $V_\alpha\in V_{\ordsucc{\alpha}}$، نتیجه
می‌شود~$V_\alpha\in V_\alpha$. پس فرض استقرایی
$V_\alpha\notin V_\alpha$ نتیجه می‌دهد
$V_{\ordsucc{\alpha}}\notin V_{\ordsucc{\alpha}}$.

اگر~$\alpha$ حدی و
$V_\alpha\in V_\alpha=\bigcup_{\beta\in\alpha}V_\beta$ باشد، برای یک
$\beta\in\alpha$ داریم~$V_\alpha\in V_\beta$؛ اما از سوی دیگر
$V_\beta\in V_\alpha$، و بنابراین با دو بار استفاده از
\olref{Valphabasicprops}، $V_\beta\in V_\beta$. پس فرض‌های استقرایی
$V_\beta\notin V_\beta$ برای همهٔ~$\beta\in\alpha$ نتیجه می‌دهند
$V_\alpha\notin V_\alpha$.
\end{proof}

\begin{cor}
برای هر دو عدد ترتیبی~$\alpha$ و~$\beta$:
$\alpha\in\beta$ اگر و تنها اگر~$V_\alpha\in V_\beta$.
\end{cor}

\begin{proof}
\Olref{Valphabasicprops} یک جهت را به دست می‌دهد. برای جهت عکس، فرض کنید
$V_\alpha\in V_\beta$. بنا بر~\olref{Valphanotref}،
$\alpha\neq\beta$؛ و~$\beta\notin\alpha$، زیرا در غیر این صورت
$V_\beta\in V_\alpha$ می‌بود و با دو بار استفاده از
\olref{Valphabasicprops} نتیجه می‌شد~$V_\beta\in V_\beta$، برخلاف
\olref{Valphanotref}. پس بنا بر سه‌شقی، $\alpha\in\beta$.
\end{proof}
\noindent
همهٔ این نتایج به ما اجازه می‌دهند هر~$V_\alpha$ را مرحلهٔ~$\alpha$ام
سلسله‌مراتب بدانیم. دلیلش چنین است.

بی‌تردید می‌توان~$V_\alpha$ها را حاصل فرایندی \emph{تکراری} دانست، زیرا
کاربرد اعداد ترتیبی مفهوم تکرار را پی می‌گیرد. افزون بر این، اگر یک مرحله
پیش از مرحله‌ای دیگر تشکیل شود، یعنی~$V_\beta\in V_\alpha$ یا به‌عبارتی
$\beta\in\alpha$، فرایند تشکیل ما \emph{انباشتی} است، زیرا
$V_\beta\subseteq V_\alpha$. سرانجام، واقعاً \emph{همهٔ} گردایه‌های ممکن
از مجموعه‌های در دسترس در مراحل پیشین را تشکیل می‌دهیم، زیرا هر مرحلهٔ
جانشین~$V_{\ordsucc{\alpha}}$ مجموعهٔ توانیِ مرحلهٔ پیشین خود، یعنی
$V_\alpha$، است.

خلاصه آنکه با~$\ZFminus$ برای بیان تصورمان از سلسله‌مراتب
انباشتی-تکراریِ مجموعه‌ها \emph{تقریباً} کار را به پایان رسانده‌ایم.
(البته هنوز باید جایگزینی را توجیه کنیم.)

\end{document}
